\documentclass[12pt,a4paper]{article}
\usepackage[utf8]{inputenc}
\usepackage[T1]{fontenc}
\usepackage[french]{babel}
\usepackage{amsmath, amssymb}
\usepackage{graphicx}
\usepackage{hyperref}
\usepackage{geometry}
\geometry{hmargin=2.5cm,vmargin=2.5cm}

\title{\textbf{Projet d'Analyse Numérique} \\ \Large Étude d'une cascade trophique}
\author{Evrard \and Romain \and Thibaud \and Zhouair}
\date{04 mai 2026} % Date de rendu prévue

\begin{document}

\maketitle
\tableofcontents
\newpage

\section{Introduction}
% C'est ta partie ! Tu dois contextualiser le sujet sur les enjeux actuels.

\section{Présentation des méthodes numériques}
% Partie de Thibaud : Les choix de méthodes numériques proposés devront être justifiés.

\section{Présentation des résultats et analyse}
% Partie de Zhouair : Graphes de l'évolution des 5 populations et Scénario Yellowstone.

\section{Conclusion}
% Ta partie : Synthèse du projet.

\newpage
\appendix
\section{Annexe : Répartition du travail et temps passé}
% Obligatoire : Indiquer comment le travail a été réparti et une estimation du temps passé par chacun.

\subsection*{Déclaration d'utilisation de l'IA}
% Le recours à l'IA doit être clairement identifié et ne doit pas dépasser 10% du rendu.

\end{document}
